\documentclass[]{article}

\usepackage{hyperref}
\usepackage{siunitx}
\usepackage{bm}
\let\vec\bm
\usepackage{xurl}
\usepackage{subcaption}
\usepackage{graphicx}   % Including figure files
\usepackage{amsmath}   % Advanced maths commands
\usepackage{amssymb}   % Extra maths symbols
\usepackage{xcolor}
%\usepackage{soul}
\newcommand{\mathcolorbox}[2]{\colorbox{#1}{$\displaystyle #2$}}
\colorlet{lgray}{gray!50!white}
\colorlet{lred}{red!50!white}
\newcommand{\bnabla}{\bm{\nabla}}
\newcommand{\bcdot}{\bm{\cdot}}

%opening
\title{Multi-scale Filtering}
\author{Lorenzo M. Perrone, Thomas Berlok, Ewald Puchwein, Christoph Pfrommer}

\begin{document}

\maketitle

\begin{abstract}
In this document we present analytical and numerical tests of multi-scale filtering of MHD turbulence, with application to galaxy cluster turbulence.
\end{abstract}

\section{Analytical Results}


\subsection{Scalar field: additive perturbation}

Let us focus in this section on the case of a simple scalar field $\rho$ of the form
%
\begin{align}
	\rho = \rho_0 + \delta \rho,
\end{align}
%
where
%
\begin{align}
	\rho_0 =\frac{\exp \left(-\frac{x^2}{2 \sigma ^2}\right)}{\sqrt{2 \pi } \sigma }
\end{align}
%
is a large-scale gaussian background profile, and
%
\begin{align}
	\delta \rho = a \cos \left(k_1 x\right)
\end{align}
%
is a small-scale, low-amplitude perturbation ($a \ll 1$). When a gaussian filter $ W_l (x)$ (where $l$ is the smoothing scale) is applied to such a scalar field, the smoothed profile $\langle \rho \rangle_l$ is 
%
\begin{align}
	\langle \rho \rangle_l \equiv \rho * W_l = \frac{e^{-\frac{x^2}{2 \left(l^2+\sigma^2\right)}}}{\sqrt{2 \pi}  \sqrt{l^2+\sigma ^2}} + \frac{e^{-\frac{1}{2} k_1^2 l^2} }{\sqrt{2 \pi }} a \cos (k_1 x),
\end{align}
%
(this is easily obtained using the convolution theorem, Eq.\eqref{app:convolution}). The reconstructed fluctuation $\delta \rho_l \equiv \rho - \langle \rho \rangle_l$ is
%
\begin{align}
	\delta \rho_l = \left( \frac{\exp \left(-\frac{x^2}{2 \sigma ^2}\right)}{\sqrt{2 \pi } \sigma } - \frac{e^{-\frac{x^2}{2 \left(l^2+\sigma^2\right)}}}{\sqrt{2 \pi}  \sqrt{l^2+\sigma ^2}} \right) + \left( 1 - \frac{e^{-\frac{1}{2} k_1^2 l^2} }{\sqrt{2 \pi }} \right) a \cos (k_1 x).
\end{align}
%
For $\sigma \gg l \gg k_1^{-1}$ (i.e. the smoothing length is larger than the wavelength of the fluctuation, but smaller than the scale-length of the background profile), $\delta \rho_l \rightarrow \delta \rho$.
%\textbf{Note that in this limit, the volume integral will yield zero!}

\subsection{Vector field: additive perturbation}

\subsubsection{Oblique wave}

Let us now consider the case of a three-dimensional vector field $\bm B (x,y,z)$  which consists of a uniform component in one direction $\bm y$ and an oblique wave $\bm k = k_x \bm x + k_y \bm y$ in the $x-y$ plane superimposed onto it:
%
\begin{align}
	\bm B (x,y,z) = \begin{pmatrix}
		B_x \\ B_y \\ B_z
	\end{pmatrix} = \begin{pmatrix}
	a \cos \left(k_x x + k_y y\right)  \\ 1 + b \cos \left(k_x x + k_y y\right) \\ 0
	\end{pmatrix},
\end{align}
%
where $b = -k_x a / k_y$ due to the incompressibility condition.
Following Hollins2022, we can write down the following relationship between the ``pseudo-energy'' density of the fluctuation ($\mu_l(B,B)$) and the energy density of the smoothed field $\langle \bm B \rangle_l$:
%
\begin{align}\label{eq:pseudo_energy}
	\langle B^2 \rangle_l = \langle \bm B \rangle_l \bcdot\langle \bm B \rangle_l + \mu_l(B,B),
\end{align} 
%
where $\langle \bm B \rangle_l = \left(\langle B_x \rangle_l , \langle B_y \rangle_l, \langle B_z \rangle_l \right)^{\text{T}}$. Applying the three-dimensional gaussian filter to $\bm B$ we get:
%
\begin{align}
	\langle \bm B \rangle_l \bcdot\langle \bm B \rangle_l =  1 + (a^2 + b^2) e^{- l^2 k^2} \cos^2 (\bm k \bcdot\bm x) + \mathcolorbox{lred}{2b e^{-\frac{1}{2} l^2 k^2} \cos (\bm k \bcdot\bm x)} ,
\end{align}
%
while 
%
\begin{align}
	\langle B^2 \rangle_l = 1 + \frac{(a^2 + b^2)}{2}  \left(e^{-2 l^2 k^2} \cos (2 \bm k \bcdot\bm x)+1\right) + \mathcolorbox{lred}{2b e^{-\frac{1}{2} l^2 k^2} \cos (\bm k \bcdot\bm x)}.
\end{align}
%
Note that the last two terms in the equations above are the same. To compute the ``pseudo-energy'' density of the fluctuations, then, we invert Eq.\eqref{eq:pseudo_energy} and obtain:
%
\begin{align}
	\mu_l(B,B) &= \langle B^2 \rangle_l - \langle \bm B \rangle_l \bcdot\langle \bm B \rangle_l  \\
	&= \frac{(a^2 + b^2)}{2}  - (a^2 + b^2) e^{- l^2 k^2} \cos^2 (\bm k \bcdot\bm x) + \\
	&+ e^{- 2 l^2 k^2} \frac{(a^2 + b^2)}{2}  \cos (2 \bm k \bcdot\bm x) .
\end{align}
%
When averaging over the entire volume, it is easy to see that the total ``pseudo-energy'' of the fluctuations will consists of an exponentially suppressed term ($\propto \exp(-k^2 l ^2)$) and a constant $(a^2 + b^2)/2$ which is exactly what we would call an ``intuitive'' measure of the turbulent energy in the domain.
If instead we decided to compute $\delta B^2_{\text{naive}} \equiv B^2 - \langle B^2 \rangle_l$ as a \textit{na\"{i}ve} measure of the turbulent energy, we would find out that after volume integration there would be no surviving term! In fact:
%
\begin{align}
	&\delta B^2_{\text{naive}} \equiv B^2 - \langle B^2 \rangle_l = a^2 \cos^2 (\bm k \bcdot\bm x) + 1 + \mathcolorbox{lgray}{2b \cos (\bm k \bcdot\bm x)} + b^2 \cos^2 (\bm k \bcdot\bm x) + \nonumber \\ 
	&- \left( 1 + \frac{(a^2 + b^2)}{2} + \mathcolorbox{lgray}{\frac{(a^2 + b^2)}{2} e^{-2 l^2 k^2} \cos (2 \bm k \bcdot\bm x)} + \mathcolorbox{lgray}{2b e^{-\frac{1}{2} l^2 k^2} \cos (\bm k \bcdot\bm x)} \right)  \nonumber \\
	&\xrightarrow{\mathcal{V}^{-1}\int_\mathcal{V} \mathrm{d} \bm x} \frac{(a^2 + b^2)}{2} + 1 - \left( 1 + \frac{(a^2 + b^2)}{2} \right) = 0,
\end{align}
%
as the terms highlighted in gray vanish after volume integration.

\subsubsection{Generic Fourier series}

Let us generalize the previous results to a real vector field with arbitrary Fourier expansion
%
\begin{align}
	\bm B (x,y,z) = \begin{pmatrix}
		B_x \\ B_y \\ B_z
	\end{pmatrix} = \begin{pmatrix}
		\sum_{\bm k} a_{\bm k} e^{i \bm k \bcdot\bm x }  \\ \sum_{\bm k} b_{\bm k} e^{i \bm k \bcdot\bm x } \\ \sum_{\bm k} c_{\bm k} e^{i \bm k \bcdot\bm x }
	\end{pmatrix},
\end{align}
%
where the wavevector $\bm k = 2 \pi \left(n_x/L_x, n_y/L_y, n_z/L_z \right)^\mathrm{T}$, and $\left(n_x, n_y, n_z \right) \in \mathbb{Z}^3$, while the coefficients $a_{\bm k}$, $b_{\bm k}$, $c_{\bm k}$ are complex Hermitian ($a_{-\bm k} = a_{\bm k}^*$), possibly related by the incompressibility condition $k_x a_{\bm k} + k_y b_{\bm k} + k_z c_{\bm k} = 0, \forall \bm k$. 
The smoothed field $\langle \bm B \rangle_l$ is (see Appendix~\ref{app:smoothing_exp}):
%
\begin{align}
	\langle \bm B \rangle_l = \begin{pmatrix}
		\langle B_x \rangle_l \\ \langle B_y \rangle_l \\ \langle B_z \rangle_l
	\end{pmatrix} = \begin{pmatrix}
		\sum_{\bm k} a_{\bm k} e^{i \bm k \bcdot\bm x } e^{-\frac{1}{2} k^2 l^2}   \\ \sum_{\bm k} b_{\bm k} e^{i \bm k \bcdot\bm x } e^{-\frac{1}{2} k^2 l^2}  \\ \sum_{\bm k} c_{\bm k} e^{i \bm k \bcdot\bm x } e^{-\frac{1}{2} k^2 l^2} 
	\end{pmatrix},
\end{align}
%
where $k^2 = \bm k \bcdot\bm k$, so that its squared norm $\langle \bm B \rangle_l^* \bcdot\langle \bm B \rangle_l$ yields (showing for simplicity only the terms relative to $B_x$):
%
\begin{align}
	&\langle B_x \rangle_l^* \bcdot\langle B_x \rangle_l  = \sum_{\bm k} |a_{\bm k}|^2 e^{- k^2 l^2} + 2 \sum_{\bm k' \neq \bm k} \Re \lbrace a_{\bm k}^* a_{\bm k'} e^{i (\bm k' - \bm k) \bcdot\bm x } e^{-\frac{1}{2} (k'^2 + k^2) l^2} \rbrace,
\end{align}
%
and similarly for the $B_y$ and $B_z$ terms.
We compute the smoothed energy $\langle \bm B^* \bcdot\bm B \rangle_l = \langle B_x^* B_x + B_y^* B_y + B_z^* B_z \rangle_l$ component-wise. For $\langle B_x^* B_x \rangle_l$ we have:
%
\begin{align}
	\langle B_x^* B_x \rangle_l &= \langle \sum_{\bm k} |a_{\bm k}|^2  + 2 \sum_{\bm k' \neq \bm k} \Re \lbrace a_{\bm k}^* a_{\bm k'} e^{i (\bm k' - \bm k) \bcdot\bm x } \rbrace  \rangle_l = \nonumber \\
	&= \sum_{\bm k} |a_{\bm k}|^2 + 2 \sum_{\bm k' \neq \bm k} \Re \lbrace a_{\bm k}^* a_{\bm k'} e^{i (\bm k' - \bm k) \bcdot\bm x } e^{-\frac{1}{2} (\bm k' - \bm k)^2 l^2} \rbrace,
\end{align}
%
and similarly for the $B_y^* B_y$ and $B_z^* B_z$ terms.
%We rewrite the term $e^{-\frac{1}{2} (\bm k_i + \bm k_j)^2 l^2}$ as $e^{-\frac{1}{2} (k_i^2 + k_j^2) l^2} e^{-(\bm k_i \bcdot\bm k_j)l^2}$, and similarly for $e^{-\frac{1}{2} (\bm k_i - \bm k_j)^2 l^2} = e^{-\frac{1}{2} (k_i^2 + k_j^2) l^2} e^{(\bm k_i \bcdot\bm k_j)l^2}$. We see that some simplifications occur when subtracting the \colorbox{green}{green} term and the \colorbox{yellow}{yellow} term. We can in fact rewrite the $\cos$ product in the \colorbox{yellow}{yellow} term using the multiplication formulae as:
%%
%\begin{align}
%	&\mathcolorbox{yellow}{\sum_{i>j} 2 a_i a_j e^{-\frac{1}{2} (k_i^2 + k_j^2) l^2} \cos \left( \bm k_i \bcdot\bm x  + \alpha_i \right) \cos \left( \bm k_j \bcdot\bm x  + \alpha_j \right)} = \nonumber \\
%	&\sum_{i>j} a_i a_j e^{-\frac{1}{2} (k_i^2 + k_j^2) l^2} \left[  \cos \left( (\bm k_i + \bm k_j) \bcdot\bm x  + \alpha_i + \alpha_j \right) +  \cos \left( (\bm k_i - \bm k_j) \bcdot\bm x  + \alpha_i - \alpha_j \right) \right].
%\end{align}
%
For simplicity we also split the ``pseudo-energy'' of the fluctuations $\mu_l(B,B) = \langle \bm B^* \bcdot\bm B \rangle_l - \langle \bm B \rangle_l^* \bcdot\langle \bm B \rangle_l$ between the $\mu_l(B_x,B_x)$, $\mu_l(B_y,B_y)$ and $\mu_l(B_z,B_z)$ components:
%
\begin{align}\label{eq:pseudo-energy-Bx}
	&\mu_l(B_x,B_x) = \langle B_x^* B_x \rangle_l - \langle B_x \rangle_l^* \bcdot\langle B_x \rangle_l = \mathcolorbox{yellow}{\sum_{\bm k} |a_{\bm k}|^2 \left( 1 -  e^{- k^2 l^2} \right)} \nonumber \\
	& + \mathcolorbox{green}{2 \sum_{\bm k' \neq \bm k} \Re \lbrace a_{\bm k}^* a_{\bm k'} e^{i (\bm k' - \bm k) \bcdot\bm x } \rbrace \left[ e^{-\frac{1}{2} (\bm k' - \bm k)^2 l^2} - e^{-\frac{1}{2} (k'^2 + k^2) l^2}\right]} .
\end{align}
%
In Eq.~\eqref{eq:pseudo-energy-Bx} we have a \colorbox{yellow}{uniform} term and \colorbox{green}{spatially-varying} term. The uniform term is directly related to the amount of energy in the small scales. In fact, assuming a given spectral distribution of energy in $\bm k$ space (e.g., a power-law), 
%is rather cumbersome, but we can immediately see that, 
for all the $\bm k$ such that the filter length is much larger than than  the wavelength of the mode ($l \gg 1/k$), the series in the \colorbox{yellow}{first term} will approximately be
%the only terms that do not vanish after integrating over the volume are
%
\begin{align}
	\mathcolorbox{yellow}{\mu_l(B,B)} \approx \sum_{|\bm k| \gg 1/l} \left( |a_{\bm k}|^2 + |b_{\bm k}|^2 + |c_{\bm k}|^2 \right),
\end{align}
%
which corresponds to our intuitive definition of the energy contained in the small-scale modes. The \colorbox{green}{spatially-varying} term vanishes exactly when integrated over the entire volume. When integration is carried out on a smaller sub-volume it is not guaranteed to be exactly zero, but should be sufficiently small provided that the extent of the region of integration is $\gg 1/| \bm k' - \bm k|$. \textbf{Note}: this condition may not be fully satisfied when $\bm k \approx \bm k'$ (but the two wavevectors in the \colorbox{green}{spatially-varying} term cannot be exactly the same).
%while for the $\bm k_i, \bm k_j$ such that $l \ll 1/k_i \sim 1/k_j$, the exponentials in the \colorbox{green}{green} and \colorbox{yellow}{yellow} terms are $\approx 1$ and therefore they approximately cancel each other. The remaining terms that do not vanish after volume integration (the constant term and the cosine squared) give:
%%
%\begin{align}
%	\sum_i \frac{a_i^2}{2} \left( 1 - e^{- k_i^2 l^2} \right) \sim \mathcal{O} (k_i^2 l^2)  \approx 0
%\end{align}
%%
%and similarly for the $y,z$ components. When $1/k_i \ll l \ll 1/k_j$, i.e. $\bm k_i$ is small-scale (fluctuation), and $\bm k_j$ is large-scale (bulk). In this case $k_j$ can be neglected in all exponentials and the \colorbox{green}{green} and \colorbox{yellow}{yellow} cross products also approximately cancel each other.
%In conclusion: when a gaussian filter of size $l$ is applied to a generic vector field, only the wavenumbers $1/k_i \ll l$ will contribute to the ``pseudo-energy'' of the fluctuations, the large-scale modes are suppressed


\subsubsection*{Inadequacy of the \textit{na\"{i}ve} estimate of the turbulent energy}

In much the same way as for the oblique mode on top of a uniform background, also in the more general case of a Kolmogorov power spectrum the \textit{na\"{i}ve} estimate of the turbulent magnetic energy $\delta B^2_{\text{naive}} \equiv \bm B^* \bcdot\bm B - \langle \bm B^* \bcdot\bm B \rangle_l  $ fails, as it averages to zero when integrated over the volume. In fact:
%
\begin{align}
	&\delta B^2_{\text{naive}} \equiv \bm B^* \bcdot\bm B - \langle \bm B^* \bcdot\bm B \rangle_l =  B_x^*  B_x - \langle B_x^*  B_x \rangle_l + \ldots \text{($B_y$ and $B_z$ terms)} \nonumber \\
	& = \sum_{\bm k} |a_{\bm k}|^2  + 2 \sum_{\bm k' \neq \bm k} \mathcolorbox{lgray}{\Re \lbrace a_{\bm k}^* a_{\bm k'} e^{i (\bm k' - \bm k) \bcdot\bm x } \rbrace} \nonumber \\
	& - \left[ \sum_{\bm k} |a_{\bm k}|^2  + 2 \sum_{\bm k' \neq \bm k} \mathcolorbox{lgray}{\Re \lbrace a_{\bm k}^* a_{\bm k'} e^{i (\bm k' - \bm k) \bcdot\bm x } \rbrace} e^{-\frac{1}{2} (\bm k' - \bm k)^2 l^2} \right] + \nonumber \\
	& + \ldots \text{($B_y$ and $B_z$ terms)} \xrightarrow{\mathcal{V}^{-1}\int_\mathcal{V} \mathrm{d} \bm x}  0,
\end{align}
%
as the terms highlighted in gray vanish exactly after volume integration. \textbf{Note}: this was also verified numerically, where due to roundoff and integration accuracy it can happen that the volume average of $\delta B^2_{\text{naive}}$ is negative.


\subsection{Convergence of iterative filter}

In this section we study the behaviour of the iterative multi-scale filter \textit{\'{a} la} Vazza in the presence of a broad spectral energy distribution (e.g., a power-law). The aim of this section is to write down sufficient conditions for convergence, and which form of the convergence criterion is more appropriate. We begin with a scalar field and then move on to a vector field.

\subsubsection{Scalar field}

As before, let us consider a scalar field $\rho (\bm x)$ with arbitrary Fourier expansion
%
\begin{align}\label{eq:def_rho_scalar}
	\rho (\bm x) = \sum_{\bm k} \hat{\rho}_{\bm k} e^{i \bm k \bcdot\bm x } ,
\end{align}
%
where the wavevector $\bm k = 2 \pi \left(n_x/L_x, n_y/L_y, n_z/L_z \right)^\mathrm{T}$, and $\left(n_x, n_y, n_z \right) \in \mathbb{Z}^3$, and with complex Hermitian coefficients ($\hat{\rho}_{-\bm k} = \hat{\rho}_{\bm k}^*$). 
The smoothed field $\langle \rho \rangle_l$ is (see Appendix~\ref{app:smoothing_exp}):
%
\begin{align}
	\langle \rho \rangle_l = \sum_{\bm k} \hat{\rho}_{\bm k} e^{i \bm k \bcdot\bm x } e^{-\frac{1}{2} k^2 l^2},
\end{align}
%
where $k^2 = \bm k \bcdot\bm k$. The small-scale component is then:
%
\begin{align}
	\delta \rho_l \equiv \rho - \langle \rho \rangle_l = \sum_{\bm k} \hat{\rho}_{\bm k} e^{i \bm k \bcdot\bm x } \left(1 - e^{-\frac{1}{2} k^2 l^2} \right).
\end{align}
%
In the original iterative procedure introduced by Vazza, the filter length $l$ is gradually increased in discrete steps $\Delta l$ until the following criterion is satisfied for each point in space:
%
\begin{align}\label{eq:Vazza_convergence}
	\lvert \delta \rho_{l+\Delta l} - \delta \rho_l \rvert \leq \epsilon \lvert  \delta \rho_l \rvert,
\end{align}
%
with $\epsilon$ a small number. If consider $\delta \rho_l$ as a function of the continuous variable $l$ and we Taylor-expand Eq.~\eqref{eq:Vazza_convergence} around $\Delta l = 0$, the above criterion for convergence can be rewritten as:
%
\begin{align}
	\Delta l \left\lvert \frac{\mathrm{d} \delta \rho_l }{\mathrm{d} l} \right\rvert \leq \epsilon \lvert  \delta \rho_l \rvert.
\end{align}
%
If we further assume that $\epsilon \sim \mathcal{O} (\Delta l / l)$, then the convergence criterion takes the simple form
%
\begin{align}\label{eq:log_convergence}
	\left\lvert\frac{\mathrm{d} \ln \delta \rho_l}{\mathrm{d} \ln l} \right\rvert \leq \eta, \hspace{2em} \text{for} \; \delta \rho_l \neq 0,
\end{align}
%
where $\eta \sim \mathcal{O}(1)$. The physical meaning of Eq.~\eqref{eq:log_convergence} is that in order for the iterative filter to converge, the relative variation in absolute value of $\delta \rho_l$ at the $n$-th iteration must be less than the relative change in the filtering radius $l$, i.e., the ``new'' information seen by the smoothing kernel does not substantially change the previous estimate of the fluctuating component $\delta \rho_l$. Note the interesting property that for a Gaussian kernel $W_l$ of width $l$ the derivative of $\delta \rho_l$ corresponds to a Mexican hat filter applied to $\rho(\bm x)$. In fact:
%
\begin{align}\label{eq:delta_density_mexican}
	\frac{\mathrm{d} \delta \rho_l}{\mathrm{d} l} &= \frac{\mathrm{d} }{\mathrm{d} l} \left( \rho - \langle \rho \rangle_l \right) = -   \int_{-\infty}^{\infty} \rho(\bm x') \frac{\mathrm{d} }{\mathrm{d} l} W_l(\bm x-\bm x')\, \mathrm{d}\bm x' \nonumber \\
	&= \frac{1}{l} \int_{-\infty}^{\infty} \rho(\bm x') M_l(\bm x-\bm x')\, \mathrm{d}\bm x',
\end{align}
%
see Appendix~\ref{app:gaussian_mexican}. The main issue with the criterion in Eq.~\eqref{eq:Vazza_convergence} or Eq.~\eqref{eq:log_convergence} is that there may be points in the domain where $\delta \rho_l = 0$ for some $l$. While in general $\delta \rho_{l'} \neq 0$ for $l' \neq l$ this might still pose a problem for the convergence of the algorithm. One option is to replace Eq.~\eqref{eq:log_convergence} with
%
\begin{align}\label{eq:log_convergence_mod1}
	\left\lvert\frac{\mathrm{d} \delta \rho_l}{\mathrm{d} \ln l} \right\rvert \leq \eta \; \text{max} \lbrace \left\lvert \rho - \text{sup}_{\bm x} \rho \right\rvert, \left\lvert \rho - \text{inf}_{\bm x} \rho \right\rvert  \rbrace,
\end{align}
%
where now the right-hand side is guaranteed to be strictly positive (unless $\rho (\bm x) = \text{const.}$, in what case the left-hand side is also zero).

Using the result in Eq.~\eqref{eq:delta_density_mexican} and the criterion in Eq.~\eqref{eq:log_convergence_mod1} we can determine a sufficient condition for convergence as follows:
%
\begin{align}\label{eq:inequality_series}
	\left\lvert\frac{\mathrm{d} \delta \rho_l}{\mathrm{d} \ln l} \right\rvert &= \left\lvert \int_{-\infty}^{\infty} \rho(\bm x') M_l(\bm x-\bm x')\, \mathrm{d}\bm x' \right\rvert = \left\lvert \sum_{\bm k} \hat{\rho}_{\bm k} k^2 l^2 e^{-\frac{1}{2} k^2 l^2} e^{i \bm k \bcdot\bm x} \right\rvert \nonumber \\
	&\leq  \sum_{\bm k} \left\lvert  \hat{\rho}_{\bm k} k^2 l^2 e^{-\frac{1}{2} k^2 l^2} e^{i \bm k \bcdot\bm x} \right\rvert = \sum_{\bm k} \left\lvert  \hat{\rho}_{\bm k} \right\rvert k^2 l^2 e^{-\frac{1}{2} k^2 l^2} ,
\end{align}
%
where in going from the convolution integral to the Fourier representation we used the Fourier expansion of $\rho (\bm x)$ in Eq.~\eqref{eq:def_rho_scalar} and the expression for the smoothing of a complex exponential with the Mexican hat filter in Eq.~\eqref{eq:mex_smoothing_Fourier}. If we now assume that the density field has a power-law power spectrum with slope $-\alpha$ (i.e. $\left\lvert\hat{\rho}(\bm k)\right\rvert \propto k^{-\alpha / 2 - 1}$, see Eq.~\eqref{eq:power_law_fourier_coeff}), we can approximate the series in Eq.~\eqref{eq:inequality_series} by the integral over the continuous variable $\mathrm{d} \bm k$
%
\begin{align}\label{eq:iterative_crit_integral}
	&\sum_{\bm k} k^{-\alpha / 2 - 1} k^2 l^2 e^{-\frac{1}{2} k^2 l^2} \approx \frac{1}{(\Delta k)^3} \int k^{-\alpha / 2 - 1} k^2 l^2 e^{-\frac{1}{2} k^2 l^2} \mathrm{d} \bm k \nonumber \\
	&= \frac{L_x L_y L_z}{(2 \pi)^3} \int_0^{\infty} k^{-\alpha / 2 - 1} k^2 l^2 e^{-\frac{1}{2} k^2 l^2} 4 \pi k^2 \mathrm{d} k \nonumber \\
	&= \frac{L_x L_y L_z}{2 \pi^2} 2^{1-\frac{\alpha }{4}} \Gamma \left(2-\frac{\alpha }{4}\right) l^{\frac{\alpha }{2}-2} \hspace{2em}\text{ if }\alpha <8,
\end{align}
%
where $\Gamma(z)$ is the gamma function. The expression in Eq.~\eqref{eq:iterative_crit_integral} is a decreasing function of $l$ provided that $\alpha < 4$ (for comparison, $\alpha = 5/3 \simeq 1.67$ for Kolmogorov turbulence and $\alpha = 3$ for Burger turbulence), which means that convergence of the iterative scheme using criterion in Eq.~\eqref{eq:log_convergence_mod1} is guaranteed for large enough $l$.

\subsection{Vector field: multiplicative perturbation}

Let us now consider a real vector field composed of the product of a large-scale smooth component $ g(\bm x)$ (e.g., a Gaussian), and a power-law that extends from intermediate- to small-scales, which we identify with the ``perturbation'' $\delta \bm B(\bm x)$.
%
\begin{align}\label{eq:mag_field_multiplicative}
	\bm B (x,y,z) = B_0  g(\bm x) \left[ \bm 1 + \delta \bm B(\bm x) \right].
\end{align}
%
The functions $g (\bm x)$ and $\delta \bm B(\bm x)$ can be Fourier-expanded as
%
\begin{align}
	g (\bm x) = \sum_{\bm k}  \hat{g}_{\bm k} e^{i \bm k \bcdot \bm x },     \hspace{2em} \delta \bm B (\bm x) = \sum_{\bm k}  \hat{ \bm b}_{\bm k} e^{i \bm k \bcdot \bm x } - \bm 1,
\end{align}
%
where $\hat{g}_{\bm k}$ and $\hat{\bm b}_{\bm k}$ are the complex Hermitian vector Fourier transform ($\hat{g}_{-\bm k} = \hat{g}_{\bm k}^*$, $\hat{\bm  b}_{-\bm k} = \hat{\bm b}_{\bm k}^*$) corresponding to wavevector $\bm k = 2 \pi \left(n_x/L_x, n_y/L_y, n_z/L_z \right)^\mathrm{T}$, with $\left(n_x, n_y, n_z \right) \in \mathbb{Z}^3$. For a Gaussian smooth component
%
\begin{align}
	g(\bm x) = \frac{\exp \left(-\frac{\bm x^2}{2 \sigma ^2}\right)}{\sqrt{2 \pi } \sigma }
\end{align}
%
the Fourier coefficients are given by Eq.~\eqref{eq:fourier_coeff_gaussian}. The magnetic field in Eq.~\eqref{eq:mag_field_multiplicative} can be always rewritten as a single Fourier series
%
\begin{align}
	\bm B (\bm x) = B_0 \sum_{\bm k}  \hat{ \bm \beta}_{\bm k} e^{i \bm k \bcdot \bm x },
\end{align}
%
where the Fourier coefficients  $\hat{ \bm \beta}_{\bm k}$ are related by the incompressibility condition $k_x \hat{\beta}_{x,\bm k} + k_y \hat{\beta}_{y,\bm k} + k_z \hat{\beta}_{z,\bm k} = 0, \forall \bm k$, and are given by
%
\begin{align}
	\hat{\bm \beta}_{\bm k} = \sum_{\bm k'}  \hat{g}_{\bm k - \bm k'} \hat{\bm b}_{\bm k'}. 
\end{align}

\newpage

\section{Numerical Tests}

In this section we show some numerical tests to validate and test the applicability of the previous results.


\subsection{Uniform background + single frequency fluctuation}

We initialize a magnetic field which consists of a uniform term plus a sinusoidal fluctuation at wavenumber $k_y = 10 \times 2 \pi / L_y$ in one of the components:
%
\begin{align}
	\bm B (\bm x) = a \cos (k_y y) \bm e_x + B_0 \bm e_y,
\end{align}
%
where $B_0 = 1 \si{G}$ and $a = 0.3\si{G}$. 

 
\begin{figure}
	\centering
	\includegraphics[width=1.2\linewidth]{Figures/task2-plots/magnetic_energy_fixed_filter_singlewave.pdf}
	\caption{Magnetic energy and pseudo-energy of the fluctuations.
	}
	\label{fig:magnetic_energy_filtered_singlewave}
\end{figure}

\begin{figure}
	\centering
	\includegraphics[width=0.9\linewidth]{Figures/task2-plots/magnetic_energy_vs_filter_length_zoom8.pdf}
	\caption{Pseudo-energy of a single wave as a function of the filtering length $l$. The solid red line is the ``energy'' of the \textit{smoothed} magnetic field, while the green is the ``pseudo-energy'' of the turbulent field. Their sum gives the volume averaged of the smoothed energy (orange solid line). The dashed lines represent the analytical values of the large-scale and fluctuating energy. 
	}
	\label{fig:magnetic_energy_vs_filter_length_single_wave}
\end{figure}

\subsection{Kolmogorov synthetic magnetic turbulence}

We now consider a more realistic configuration where the magnetic field is the sum of fluctuations at many different wavenumbers which follow a power-law in spectral space (e.g. Kolmogorov). 

\subsubsection{Initialization of the synthetic turbulent field}

In three dimensions, the three-dimensional energy density $P_{\text{3D}} (\bm k) \propto |a(\bm k)|^2$, with $a(\bm k)$ the complex Fourier coefficient at wavevector $\bm k$, and the one-dimensional spectral energy density $E(k)$ are related by
%
\begin{align}
	P_{\text{3D}} (\bm k) \mathrm{d} \bm k = 4 \pi P_{\text{3D}} (\bm k) k^2 \mathrm{d} k = E(k) \mathrm{d} k .
\end{align}
%
Therefore, in order to initialize a synthetic turbulent field with given one-dimensional spectral density energy $E(k) \propto k^{-\alpha}$, the Fourier coefficients at each wavevector must satisfy
%
\begin{align}\label{eq:power_law_fourier_coeff}
	a(\bm k) \propto \sqrt{\frac{E(k)}{k^2}} e^{i\phi}  \propto k^{-\alpha /2 - 1} e^{i\phi},
\end{align}
%
where $\phi$ is a random phase.

The procedure that we adopt is as follows:
%
\begin{enumerate}
	\item we load a region of space (typically corresponding to a cube of size $L_x \times L_y \times L_z$) from an \textsc{arepo} zoom-in snapshot;
	\item we define a three-dimensional $M \times M \times M$ uniform grid of wavevectors where the $x$, $y$, and $z$ components are given respectively by 
	%
	\begin{align}
		k_{x,y,z} = \frac{2 \pi}{L_{x,y,z}} \times 
		\begin{cases}
			\lbrace -(M-1)/2, \ldots , (M-1)/2 \rbrace ~~~ &\text{when $M$ is odd}, \\
			\lbrace -M/2, \ldots , M/2 - 1 \rbrace ~~~ &\text{when $M$ is even};
		\end{cases}
	\end{align} 
	\item we define the complex Fourier coefficients at each wavenumber and for each component of the magnetic field, and we enforce their Hermitianity ($a_{-\bm k} = a_{\bm k}^*$);
	\item we use the ``type 2'' (uniform to non-uniform) routines of \textsc{finufft} to compute the synthetic turbulent magnetic field on the Voronoi centers in the region of interest.
\end{enumerate}
%
An example of a synthetic turbulent magnetic field ($M=129$) is in the top-left panel of Fig.~\ref{fig:magnetic_energy_filtered}.
%
To compute the one-dimensional power spectrum we can simply use the complex amplitudes:
%
\begin{equation}
	E(k) dk = \sum_{k \leq \lvert \bm k \rvert < k + dk} \lvert a_{\bm k} \rvert^2 .
\end{equation}
%
A typical Kolmogorov power spectrum is in Fig.~\ref{fig:kolmogorov_power_spectrum}.

\begin{figure}
	\centering
	\includegraphics[width=0.75\linewidth]{Figures/task2-plots/2a.power_spectrum_mag_field.pdf}
	\caption{Power spectrum of magnetic fluctuations.
	}
	\label{fig:kolmogorov_power_spectrum}
\end{figure}

\subsubsection{Filtering of the magnetic energy}

In Fig.~\ref{fig:magnetic_energy_filtered} we show the result of the filtering procedure (with fixed filter length $l=25 \si{ckpc.h^{-1}}$). We compare the \textit{na\"{i}ve} estimate of the turbulent energy with the ``pseudo-energy'' introduced by Hollins.

\begin{figure}
	\centering
	\includegraphics[width=1.2\linewidth]{Figures/task2-plots/2a.magnetic_energy_fixed_filter_zoom12.pdf}
	\caption{Magnetic energy and pseudo-energy of the fluctuations.
	}
	\label{fig:magnetic_energy_filtered}
\end{figure}

\begin{figure}
	\centering
	\includegraphics[width=0.9\linewidth]{Figures/task2-plots/2a.magnetic_energy_vs_filter_length_zoom12.pdf}
	\caption{Pseudo-energy of small scales as a function of the filtering length $l$. The blue dashed line is the total energy of the synthetic turbulent field, the solid red line is the ``energy'' of the \textit{smoothed} magnetic field, while the green is the ``pseudo-energy'' of the turbulent field. Their sum gives the volume averaged of the smoothed energy (orange solid line). In black is the sum of the energy at small-scales obtained by applying a high-pass filter to the spectral coefficients. The black squares correspond instead to exponentially suppressed uniform term in Eq.~\eqref{eq:pseudo-energy-Bx}, i.e. $\sum_{\bm k} |a_{\bm k}|^2 \left( 1 -  e^{- k^2 l^2} \right)$.
	}
	\label{fig:magnetic_energy_vs_filter_length}
\end{figure}

\subsubsection{Is the ``pseudo-energy'' a good measure of the turbulent energy?}

In this Section we try to answer the question of whether $\mu_l (B,B)$ is a ``good'' measure of the turbulent energy on scales $< l$, i.e. it has a clear relationship to our intuitive notion of small-scale turbulent energy. To do so, we use the synthetic magnetic field shown in the previous section and apply our filtering machinery progressively increasing the filtering length (\textbf{note} this is not the same as the iterative version of our filtering kernel, because here at each step the filter length is the same across the entire volume).

In Fig.~\ref{fig:magnetic_energy_vs_filter_length} we plot the result of this procedure. We compare two analytic measures of the turbulent energy content at small scales: 
%
\begin{itemize}
	\item we apply a high-pass filter to the spectral coefficients, with cutoff $1/l$ and compute the corresponding energy (solid black line);
	\item we exponentially suppress the spectral coefficients (as in the \colorbox{yellow}{uniform term} in Eq.~\eqref{eq:pseudo-energy-Bx}), i.e. $\sum_{\bm k} |a_{\bm k}|^2 \left( 1 -  e^{- k^2 l^2} \right)$. This is shown in the black squares. 
\end{itemize}
%
The fact that there is very good agreement between the black squares and the green line $\mu(B,B)$ means that the spatially varying terms in Eq.~\eqref{eq:pseudo-energy-Bx} effectively average out to zero when volume integrated.


\subsubsection{Appropriateness of the ``pseudo-energy'' for smaller sub-volumes}

\begin{figure}
	\centering
	\includegraphics[width=0.9\linewidth]{Figures/task2-plots/2a.magnetic_energy_fixed_filter_regions.pdf}
	\caption{Slice of the magnetic energy with randomly selected subregions.
	}
	\label{fig:magnetic_energy_fixed_filter_regions}
\end{figure}

\begin{figure}
	\centering
	\includegraphics[width=0.9\linewidth]{Figures/task2-plots/2a.magnetic_energy_vs_filter_length_smaller_subregion_many.pdf}
	\caption{Pseudo-energy of small scales as a function of the filtering length $l$ for small subregions. The blue dashed line is the total energy of the synthetic turbulent field, the solid red line is the ``energy'' of the \textit{smoothed} magnetic field, while the green is the ``pseudo-energy'' of the turbulent field. Their sum gives the volume averaged of the smoothed energy (orange solid line). In black is the sum of the energy at small-scales obtained by applying a high-pass filter to the spectral coefficients. The black squares correspond instead to exponentially suppressed uniform term in Eq.~\eqref{eq:pseudo-energy-Bx}, i.e. $\sum_{\bm k} |a_{\bm k}|^2 \left( 1 -  e^{- k^2 l^2} \right)$.
	}
	\label{fig:magnetic_energy_vs_filter_length_smaller)subregions}
\end{figure}

\newpage

\appendix

\section{Some useful Fourier transforms}

The Fourier transform $\hat{\rho}_0$ of a gaussian $\rho_0 (x)$ (unitary convention, using angular frequency $k = 2 \pi / x$)
%
\begin{align}
\rho_0 (x) = \frac{\exp \left(-\frac{x^2}{2 \sigma ^2}\right)}{\sqrt{2 \pi } \sigma }
\end{align}
%
is the following:
%
\begin{align}\label{eq:fourier_coeff_gaussian}
	\hat{\rho}_0 (k) = \frac{e^{-\frac{1}{2} k^2 \sigma ^2}}{\sqrt{2 \pi }},
\end{align}
%
and similarly for a gaussian filter $W_l (x)$ with filtering length $l$ (replace $\sigma \rightarrow l$ in the equation above).


\section{Convolution theorem}

For two functions $\rho (x)$ and $ W_l (x)$ with Fourier transforms $\hat{\rho}$ and $ \hat{W}_l$  (unitary convention, using angular frequency $k = 2 \pi / x$), the convolution theorem states that:

\begin{align}\label{app:convolution}
	( \rho * W_l) (x) =  \int_{-\infty}^{\infty} \rho(\tau) W_l(x-\tau)\, d\tau = \sqrt{2 \pi} \mathcal{F}^{-1} \left[ \hat{\rho} \hat{W}_l \right].
\end{align}

\section{Smoothing of a complex exponential}\label{app:smoothing_exp}

The convolution theorem with a gaussian kernel applied to a complex exponential function yields:
%
\begin{align}
	\langle e^{i \bm k \bcdot\bm x} \rangle_l = e^{-\frac{1}{2} k^2 l^2} e^{i \bm k \bcdot\bm x} ,
\end{align}
%
where $k^2 = \bm k \bcdot\bm k$.

\section{Relationship between Gaussian kernel and Mexican hat filter}\label{app:gaussian_mexican}

In this section we show that the first derivative of a Gaussian kernel $W_l$,
%
\begin{align}
	W_l (\bm x) = \frac{e^{-\frac{x^2}{2 l ^2}}}{(\sqrt{2 \pi } l )^3},
\end{align}
%
with respect to the width $l$ corresponds to a Mexican hat filter. In fact:
%
\begin{align}\label{eq:derivative_gaussian}
	\frac{\mathrm{d} W_l (\bm x)}{\mathrm{d} l} = - \frac{1}{l}\left[3 - \frac{x^2}{l^2} \right]  W_l (\bm x) = l \nabla^2_{\bm x} W_l (\bm x) \propto - M_l (\bm x),
\end{align}
%
where $M_l (\bm x)$ is the Mexican hat filter, generally defined as the negative normalized second derivative of a Gaussian function\footnote{\url{https://en.wikipedia.org/wiki/Ricker_wavelet}}. Note that the equivalence in Eq.~\eqref{eq:derivative_gaussian} stems from the fact the the gaussian is the solution to the diffusion equation $\partial_t W = (1/2) \nabla^2 W$ (replace $l^2 \rightarrow t$). In this document we use the operative definition
%
\begin{align}
	M_l (\bm x) = - l \frac{\mathrm{d} W_l (\bm x)}{\mathrm{d} l}.
\end{align}
%
As a result, the smoothing of a complex exponential with the Mexican hat filter yields:
%
\begin{align}\label{eq:mex_smoothing_Fourier}
	( e^{i \bm k \bcdot\bm x} * M_l) (\bm x) = k^2 l^2 e^{-\frac{1}{2} k^2 l^2} e^{i \bm k \bcdot\bm x}.
\end{align}

\end{document}
